\section{Economic Robustness and the Cost of a Decision Certificate}
\label{sec:r8_evidence}
The preceding experiments establish the operator and transfer comparisons on their original targets. We now hold the economic decision itself fixed when comparing the upper components, and then vary the economic primitives with both adjustment regimes reoptimized. Every historical comparison remains in the preceding tables. The additional evidence neither substitutes a generic welfare tolerance for a sign margin nor removes the original region when its interpretation is challenged by a perturbed transition technology.

\subsection{The Original Region and Its Moving Boundary}
\begin{table}[htbp]
\centering
\caption{The original decision corner under wealth-lottery perturbations}\label{tab:r8_spatial}
\small
\begin{tabular}{rlrrr}
\toprule
Wealth nodes & Menu & $\Delta^0$ & $\Delta^{\rm adj}$ & Relative option \\
\midrule
49 & Common & $-5.3001\times10^{-5}$ & $3.2453\times10^{-4}$ & $3.7753\times10^{-4}$ \\
49 & Full & $-5.3001\times10^{-5}$ & $3.2453\times10^{-4}$ & $3.7753\times10^{-4}$ \\
97 & Common & $6.1263\times10^{-5}$ & $4.2641\times10^{-4}$ & $3.6515\times10^{-4}$ \\
97 & Full & $6.1263\times10^{-5}$ & $4.2641\times10^{-4}$ & $3.6515\times10^{-4}$ \\
145 & Common & $9.7943\times10^{-5}$ & $4.6009\times10^{-4}$ & $3.6215\times10^{-4}$ \\
145 & Full & $9.7943\times10^{-5}$ & $4.5954\times10^{-4}$ & $3.6160\times10^{-4}$ \\
\bottomrule
\end{tabular}
\par\vspace{3pt}\begin{minipage}{0.96\textwidth}\footnotesize All rows retain 33 preference tiers, eight dates, and $(\lambda,d)=(0.25,0.45)$. Common uses the identical 1,565-action union. Full also uses the three unchanged proposals at original nodes and their specified wealth interpolation at new nodes. Both classes and both adjustment regimes are reoptimized. The original region is not silently replaced by a favorable subset.\end{minipage}
\end{table}

Table \ref{tab:r8_spatial} reproduces the original full-menu corner and its movement under finer wealth settlement. The change from 49 to 97 nodes switches the no-adjustment ranking at this corner; the adjusted ranking remains positive. The positive relative adjustment option therefore survives a perturbation under which adjustment ceases to be necessary for positive risk. This distinction is not a failure of the finite-array inequality. It identifies the part of the economic conclusion that depends on the settlement lottery.

The common-menu and prolonged-full experiments give the same qualitative result. Their menus are not silently equated: interpolation of frozen proposals at new nodes can add controls beyond the common mesh. The common-menu arm isolates the wealth-lottery change without that additional policy-menu change. All four original corners and the center are retained for every grid and both menus, with 120 class values and selected-policy replays. At the center, the no-adjustment difference remains negative on all three wealth grids. Thus the evidence maps the movement of the boundary rather than asserting that the opposite-choice phenomenon disappears everywhere.

\begin{table}[htbp]
\centering
\caption{Local indifference brackets with the calendar held fixed}\label{tab:r8_frontiers}
\small
\begin{tabular}{rrrr}
\toprule
Wealth nodes & $\lambda$ & Adjusted contract bracket & No-adjustment bracket \\
\midrule
49 & 0 & $[0.370068, 0.370117]$ & $[0.558887, 0.558936]$ \\
49 & 0.125 & $[0.339551, 0.339600]$ & $[0.515869, 0.515918]$ \\
49 & 0.25 & $[0.309082, 0.309131]$ & $[0.472998, 0.473047]$ \\
97 & 0 & $[0.320166, 0.320215]$ & $[0.513477, 0.513525]$ \\
97 & 0.125 & $[0.288135, 0.288184]$ & $[0.467773, 0.467822]$ \\
97 & 0.25 & $[0.256055, 0.256104]$ & $[0.422266, 0.422314]$ \\
145 & 0 & $[0.302783, 0.302832]$ & $[0.497852, 0.497900]$ \\
145 & 0.125 & $[0.270313, 0.270361]$ & $[0.451367, 0.451416]$ \\
145 & 0.25 & $[0.237646, 0.237695]$ & $[0.405078, 0.405127]$ \\
\bottomrule
\end{tabular}
\par\vspace{3pt}\begin{minipage}{0.96\textwidth}\footnotesize All brackets use the full prolonged menu. Bisection begins on $[0.2,0.6]$ and retains every evaluated difference. Opposite signs and continuity locate a crossing; global uniqueness is not inferred. Changing the wealth lottery shifts both boundaries without changing the one-eighth-year commitment interval.\end{minipage}
\end{table}

The indifference brackets in Table \ref{tab:r8_frontiers} explain that movement on the contract scale. Both regime boundaries shift as the wealth lottery is changed, and the interval between them remains economically nonempty at the inspected mixture weights. The full solver trace records all bisection evaluations. These are local brackets, not a proof of global monotonicity of the difference of two convex value functions.

\begin{table}[htbp]
\centering
\caption{A common nested decision region across three settlement protocols}\label{tab:r8_nested}
\small
\begin{tabular}{rrrr}
\toprule
Wealth nodes & Lower $\Delta^{\rm adj}$ & Upper $\Delta^0$ & Both signs certified \\
\midrule
49 & $6.8695\times10^{-5}$ & $-2.0913\times10^{-4}$ & Yes \\
97 & $1.7639\times10^{-4}$ & $-9.4581\times10^{-5}$ & Yes \\
145 & $2.1143\times10^{-4}$ & $-5.7401\times10^{-5}$ & Yes \\
\bottomrule
\end{tabular}
\par\vspace{3pt}\begin{minipage}{0.96\textwidth}\footnotesize The same closed region is $[0,0.125]\times[0.4,0.425]$. Bounds shown use the full-menu bank and corrected chord. Count and mesh-only lower-bank checks are also retained. This is a continuum certificate in $(\lambda,d)$ for each of three specified finite protocols, not a continuum claim over every wealth grid or the limiting diffusion. The region was selected after inspecting the earlier robustness results; it was not preregistered.\end{minipage}
\end{table}

Table \ref{tab:r8_nested} supplies an additional regional statement. It uses one common nested rectangle and separately certifies its two signs under each of the three full-menu finite protocols. The old rectangle and its failed corner are retained in Table \ref{tab:r8_spatial}; the nested region is not presented as an unannounced replacement. Its validity over intermediate contracts and mixture probabilities comes from tensor coefficient bounds, not from testing a few parameter points. The discrete set of settlement protocols is itself not promoted to a continuum of all possible discretizations.

\subsection{Which Primitive Supports the Dynamic Option?}
\begin{table}[htbp]
\centering
\caption{Dynamic mechanisms at the original region center}\label{tab:r8_mechanisms}
\small
\begin{tabular}{lrrr}
\toprule
Intervention & $10^4\Delta^{\rm adj}$ & $10^4\Delta^0$ & $10^4$ relative option \\
\midrule
Baseline & 1.9674 & -2.0933 & 4.0607 \\
Tilt $a=-0.25$ & 1.5564 & -1.3521 & 2.9085 \\
Tilt $a=0.25$ & 2.3858 & -2.8501 & 5.2359 \\
Tilt $a=1$ & 3.8022 & -5.1768 & 8.9790 \\
Normalization $s=0$ & 9.1981 & 9.1981 & 0.0000 \\
Normalization $s=0.5$ & 11.6551 & 10.5243 & 1.1309 \\
Normalization $s=1.5$ & -8.9524 & -15.8240 & 6.8716 \\
Liquidation $z=0$ & -20.0244 & -24.7728 & 4.7484 \\
Liquidation $z=2$ & 21.1887 & 17.8547 & 3.3340 \\
Zero-covariance law & 4.1345 & 0.9752 & 3.1593 \\
\bottomrule
\end{tabular}
\par\vspace{3pt}\begin{minipage}{0.96\textwidth}\footnotesize Except for the named intervention, parameters are the original $(\lambda,d,k)=(0.125,0.425,2)$. Every entry reoptimizes both risk classes in both adjustment regimes. The normalization multiplier changes $b_0(u)=1/(1-u)$, not marginal consumption utility. Values shown as zero at $s=0$ are numerically negligible, not an asserted exact zero theorem. These model counterfactuals are not empirically identified causal effects.\end{minipage}
\end{table}

The dynamic counterfactuals distinguish the reference-state payoff from ordinary risk aversion. At the center, the baseline relative adjustment option is approximately $4.061\times10^{-4}$. Removing $b_0(u)$ while preserving marginal consumption utility makes this difference numerically negligible and leaves positive risk in both regimes. Increasing its multiplier to $1.5$ instead gives a larger positive relative option while making both risk rankings nonpositive. The cardinal value of preference states is therefore a principal economic primitive of the dynamic example, not an innocuous normalization.

The common tilt illustrates why the two-stage result is not a calibration of the full economy. A tilt of one reverses the relative option in the two-stage example, whereas in the dynamic economy it increases the positive relative option at the inspected center. In the latter experiment the no-adjustment value changes because preferences continue to fluctuate, and consumption, continuation, and liquidation opportunities are endogenous. The uniform primitive theorem remains valid on its stated family; the dynamic calculation supplies the separate evidence needed for this economy.

Liquidation and covariance also move the choice boundary. Eliminating the liquidation payoff makes both rankings nonpositive while leaving a positive relative adjustment option. Doubling that payoff makes both rankings positive. The zero-covariance law gives positive risk in both regimes. These interventions retain the controls, deliberate-adjustment restriction, and positive-duration decision definition. They reveal interactions in the full model rather than crediting every change in risk choice to preference formation alone.

\begin{table}[htbp]
\centering
\caption{Optimized-moment accounting for the relative adjustment option}\label{tab:r8_features}
\small
\begin{tabular}{lr}
\toprule
Component & Contribution in baseline utility units \\
\midrule
Consumption exposure & $-5.1603\times10^{-6}$ \\
Cardinal preference-state payoff & $5.0928\times10^{-4}$ \\
Tilt exposure at baseline & $0$ \\
Operating duration & $-1.3726\times10^{-5}$ \\
Adjustment effort & $-4.9862\times10^{-5}$ \\
Liquidation payoff & $-3.4460\times10^{-5}$ \\
\bottomrule
\end{tabular}
\par\vspace{3pt}\begin{minipage}{0.96\textwidth}\footnotesize The sum is the baseline relative option. Each moment is evaluated at its own class--regime optimizer. This accounting does not assign invariant causal shares; policy changes and primitive interactions are reported separately through reoptimization.\end{minipage}
\end{table}

The optimized-moment account in Table \ref{tab:r8_features} puts the baseline magnitudes in common cardinal units. The preference-state payoff contributes about $5.093\times10^{-4}$ to the relative option, partly offset by consumption exposure, duration, effort, and liquidation terms. This is an exact accounting across four separately optimized policies, not a causal identification theorem. The deposited counterfactuals additionally evaluate the original policies under each new primitive and report their reoptimization gains. They make clear which changes are valuation of a fixed exposure and which require altered behavior.

\subsection{Attributable Computation at the Economic Accuracy Target}
\begin{table}[htbp]
\centering
\caption{A common-target economic decision certificate}\label{tab:r8_contest}
\small
\begin{tabular}{llrrrr}
\toprule
Upper & Lower bank & Lower $\Delta^{\rm adj}$ & Upper $\Delta^0$ & Upper sec. & Total sec. \\
\midrule
Chord & Full & $6.8695\times10^{-5}$ & $-5.2801\times10^{-5}$ & 2.02 & 10.60 \\
Count & Full & $6.8695\times10^{-5}$ & $-5.2801\times10^{-5}$ & 7.89 & 16.47 \\
Chord & Mesh only & $6.8695\times10^{-5}$ & $-5.2801\times10^{-5}$ & 2.02 & 13.79 \\
\bottomrule
\end{tabular}
\par\vspace{3pt}\begin{minipage}{0.96\textwidth}\footnotesize The target is the original 1,568-action economy on $[0,0.25]\times[0.4,0.45]$. Every arm is charged for kernel and proposal loading, full-menu upper-anchor optimization, feasible-policy evaluation, the arithmetic audit, and independent replay. The mesh-only bank additionally pays for its own lower-policy construction. Common setup is measured once and charged identically; these are serial single-pass measurements, not independent-machine medians. The per-class arithmetic allowance is $10^{-7}$.\end{minipage}
\end{table}

\begin{table}[tbp]
\centering\small
\caption{Initial Optimizers Beside the Original Decision Certificate}
\label{tab:r9_initial}
\begin{tabular}{llrrrr}
\toprule
Regime & Class & Value & $c$ & $\theta$ & $\pi$ \\
\midrule
Adjustment & $+$ & -0.759029726 & 0.8 & 0.2 & 0.8 \\
Adjustment & $-$ & -0.759226466 & 0.8 & 0.2 & -0.5 \\
No adjustment & $+$ & -0.794487049 & 0.8 & 0.0 & 0.8 \\
No adjustment & $-$ & -0.794277722 & 0.8 & 0.0 & -0.5 \\
\bottomrule
\end{tabular}
\begin{minipage}{0.98\textwidth}\footnotesize\medskip
Original center $(\lambda,d)=(0.125,0.425)$, noncancellable mandate, complete finite menu. Both the long permission $0.8$ and the short permission $0.5$ bind. Consumption is also at its upper limit; deliberate adjustment is at its upper limit when permitted. The nonpositive position is risky, and the no-adjustment regime retains stochastic preferences.
\end{minipage}
\end{table}

Table \ref{tab:r8_contest} compares streaming chord and count implementations on the original 1,568-action target and the original rectangle. Both use the same full-menu upper anchors and independently evaluated class policies. Both certify the two risk rankings. The count coefficients are weakly tighter, but the chord requires less upper construction work; in this experiment its tighter competitor does not change the accepted economic conclusion. The total includes model construction, frozen-proposal loading, endpoint optimization, lower evaluation, the arithmetic audit, and independent selected-control replay. The gain is therefore attached to the economic decision tolerance rather than an unrelated $10^{-3}$ welfare target.

The mesh-only lower-bank arm is evaluated against the same full-menu upper target. It also certifies both signs. Its separately charged construction is a diagnostic of attribution, not a claim to be the most efficient non-neural solver. No network is trained for these new arms: the three historical proposals define part of their common, already deposited target. Their loading and extension are charged in construction; historical network training is a sunk input reported in the retained earlier experiments, not silently relabeled as a new training speedup. Adding an identical historical construction charge to all arms changes total percentages, but not the measured difference between these upper components.

Streaming addresses the lower-bank storage issue directly. For the baseline focal-state task, sixteen policies require 827,904 bytes of control indices, and their retained intercept--duration coefficients require 2,304 bytes. The kernel and reward payload is 775,469,952 bytes. These payloads exclude common temporary arrays and process overhead; the evidence records process peak resident memory separately, without attributing a shared process peak to one upper arm. The largest memory object remains the transition representation, not the upper coefficient correction. The streamed focal certificate does not claim to preserve simultaneous access to all dates' and states' lower coefficients at the same storage cost.

The new dense audit checks 60 models, 240 class recursions, and 720 parameter-point comparisons, including first-action and no-adjustment restrictions. Full-model policies are also replayed through the original selected-control transition formula. All executed values, upper and lower coefficients, source hashes, and stage timings accompany the manuscript. The bounds concern the specified finite protocols and their documented arithmetic model. No empirical replay discrepancy is used as a substitute for the additional target-error control in \eqref{eq:r8_target_budget}.
